\documentclass[11pt,a4paper]{article}
\newcommand{\fichetype}{Exercices}
% Thème A4 commun, librement modifiable.
\usepackage[T1]{fontenc}
\usepackage[utf8]{inputenc}
\usepackage[french]{babel}
\usepackage{lmodern,amsmath,amssymb,mathtools,tikz,pgfplots}
\usepackage[margin=2.2cm,headheight=15pt]{geometry}
\usepackage{xcolor,fancyhdr,enumitem,hyperref,needspace}
\usetikzlibrary{arrows.meta,positioning,calc}
\pgfplotsset{compat=1.18,width=11cm,height=6cm}
\definecolor{Ink}{HTML}{202124}
\definecolor{Blue}{HTML}{356B9E}
\definecolor{Muted}{HTML}{737579}
\color{Ink}
\renewcommand{\familydefault}{\sfdefault}
\setlength{\parindent}{0pt}
\setlength{\parskip}{6pt}
\setlist{itemsep=3pt,topsep=4pt}
\pagestyle{fancy}\fancyhf{}
\fancyhead[L]{\small\color{Muted}MATHS LYCÉE}
\fancyhead[R]{\small\color{Muted}2026--2027}
\fancyfoot[L]{\small\color{Muted}\fichetype}
\fancyfoot[R]{\small\thepage}
\renewcommand{\headrulewidth}{0.3pt}
\hypersetup{colorlinks=true,urlcolor=Blue,linkcolor=Blue}
\newcommand{\R}{\mathbb{R}}
\newcommand{\N}{\mathbb{N}}
\newcommand{\E}{\mathbb{E}}
\newcommand{\V}{\mathrm{V}}
\newcommand{\exercice}[3]{\Needspace{6\baselineskip}\vspace{0.5cm}{\color{Blue}\small\textbf{EXERCICE #1}\hfill #2}\par{\large\bfseries #3}\par}

\begin{document}
{\small\color{Blue}PREMIÈRE SPÉCIALITÉ}\par{\LARGE\bfseries Suites : construire, représenter et conjecturer}\par{\large Exercices}\par\vspace{0.2cm}\hrule\vspace{0.4cm}
Les exercices progressent de l’application directe vers la recherche. Détailler les calculs et justifier les réponses.
\par\noindent\begin{minipage}{\linewidth}\exercice{1}{Application}{Rangs et valeurs}
On définit $u_n=2n^2-3$ pour $n\ge0$. Calculer $u_0,u_1,u_4$, écrire $u_{n+1}$ et dire quel est le quatrième terme.
\end{minipage}\par
\par\noindent\begin{minipage}{\linewidth}\exercice{2}{Application}{Calcul de proche en proche}
$u_0=3$ et $u_{n+1}=u_n^2-2$. Calculer $u_1,u_2,u_3$. Pourquoi ne peut-on pas remplacer directement $u_n$ par $n$ ?
\end{minipage}\par
\par\noindent\begin{minipage}{\linewidth}\exercice{3}{Entraînement}{Une frise de triangles}
Une frise de $n\ge1$ triangles utilise trois allumettes pour le premier triangle et deux de plus pour chaque triangle supplémentaire. Donner une récurrence puis une formule explicite pour le nombre $a_n$ d’allumettes. Calculer $a_{50}$.
\end{minipage}\par
\par\noindent\begin{minipage}{\linewidth}\exercice{4}{Entraînement}{Étudier les variations}
Étudier les variations de $u_n=n^2-5n$ pour $n\ge0$. Préciser les éventuels termes consécutifs égaux.
\end{minipage}\par
\par\noindent\begin{minipage}{\linewidth}\exercice{5}{Entraînement}{Le piège du quotient}
Pour $u_n=-2^n$, calculer $u_{n+1}/u_n$, puis $u_{n+1}-u_n$. La suite est-elle croissante ?
\end{minipage}\par
\par\noindent\begin{minipage}{\linewidth}\exercice{6}{Entraînement}{Conjectures de limites}
Conjecturer et expliquer le comportement de $a_n=5-1/(n+1)$, $b_n=3n+2$, $c_n=(-1)^n$ et $d_n=(-1)^n/(n+1)$.
\end{minipage}\par
\par\noindent\begin{minipage}{\linewidth}\exercice{7}{Entraînement}{Seuil de production}
Une machine produit $u_n=40+15n$ pièces au jour de rang $n\ge0$. Écrire un algorithme donnant le premier rang pour lequel $u_n\ge200$, puis le déterminer.
\end{minipage}\par
\par\noindent\begin{minipage}{\linewidth}\exercice{8}{Approfondissement}{Vérifier une formule proposée}
On propose $u_n=4\times3^n-2$ pour la suite $u_0=2$, $u_{n+1}=3u_n+4$. Vérifier l’initialisation et la règle. Calculer alors $u_6$.
\end{minipage}\par
\end{document}
